\scanpage{26}%
% archive sheet number: 25
poser de plus d'un \uline{foncteur fibre}
\[
\unsure{F|\mathbb{F}}:\operatorname{Rev}^{\underline{R}}(S)\qquad(\mathrm{Ens\ finis}),
\]
i.e. d'un foncteur exact\inserted{(conditions G4,G5 de. V 4)}. Dans le cas de l'exemple 1, tout point géométrique de $Z$ fournit un tel foncteur fibre, et ce foncteur est conservatif (condition G6) si on suppose que toute partie ouverte et fermée de $Z$ stable par $G$ est vide ou égale à $Z$. De même, dans le cas de l'exemple 2, on obtient un \deleted{\unsure{par}} foncteur fibre en prenant un point géométrique de $Z=\lim Z_i$. Bour obtenir un foncteur fibre dans le cas plus général de l'exemple 3, on choisit un point géométrique $\qquad S$ de $S$ \deleted{tel que} qui n'appartienne pas à\deleted{\unsure{u}} ``l'ensemble de ramification''de $\underline{R}$, i.e. tel que pour tout $S'$ dans $\mathrm{Et}_S$ et tout objet $M$ de $\underline{C}_{S'}$, \unsure{le|l'} \deleted{\unsure{objet}} \inserted{revêtement à} \deleted{à} opérateurs $\underline{r}(M)$ \deleted{sur $S'$}\inserted{(groupe $G$)} de $S'$ soit principal en les points de $S'$au dessus du l'image \inserted{$\unsure{s|S}$} de $\quad$ dans $S$. Alors \deleted{pour tout revêtement étale} \deleted{$T$ \unsure{à groupe $G$} \unsure{de} \unsure{à l'ob|l'ob} d'un schéma à opérateurs $G$} on constate que pour tout objet $T$ de $\operatorname{Rev}^{\underline{R}}(S)$, $\underline{r}(T)$ est étale enr $S$ en $s$, et le foncteur $T\quad\underline{r}(T)(\quad)$ (fibre géométrique en $\quad$) est bien exact. Pour qu'il soit conservatif, il faut et il suffit que le site $\underline{C}$ (pour la topologie image inverse de celle de $\underline{\mathrm{Et}}_S$) soit connexe; \deleted{i.e. que tout crible couvrant de ce site} si on suppose, avec les notations de l'exemple 3, que pour tout $M$, définissant $(\underline{r}(M),G)$, $\underline{r}(M)/G\qquad S'$ soit un isomorphisme (ce qui sera sans doute le cas dans tous les cas intéressants), alors il revient au même de dire que $S$ est connexe. Appliquant la théorie de V 4, on trouve donc un groupe pro-fini, noté ${}^{\underline{R}}_1(S,\quad)$, \deleted{qui classifi} et une équivalence de la catégorie des ensembles finis où ce groupe opère continûment, avec la catégorie $\operatorname{Rev}^{\underline{R}}(S)$.

3. \uline{Fonctorialités pour} $S$ \inserted{\uline{et} $\underline{R}$} \uline{variables}

\unsure{Crnsidérons|Considérons} un morphisme $f:T\qquad S$ de schémas, d'où un foncteur
